% !TeX root = ../main.tex
% figures/annotated_source.tex  --  included via \input{figures/annotated_source.tex}
%
% Required in the preamble:
%   \usepackage{listings}
%   \usepackage{xcolor}
%
% NOTE: the gutter markers are typeset through listings' escape mechanism
% (escapeinside={@}{@}) rather than through `literate`, because `literate`
% inside \lstdefinestyle is unreliable across listings versions and is silently
% overridden by a global \lstset{literate=...} in some preambles. The escape
% character is @, which must therefore not occur in the listed source.

\definecolor{mkA}{HTML}{7C3AED} % [A] affected
\definecolor{mkC}{HTML}{15803D} % [C] covered
\definecolor{mkP}{HTML}{D97706} % [P] partially covered
\definecolor{mkU}{HTML}{B91C1C} % [.] uncovered
\definecolor{srcKw}{HTML}{1D4ED8}
\definecolor{srcCm}{HTML}{6B7280}
\definecolor{srcSt}{HTML}{0F766E}

% coverage-gutter markers; \annB is an empty, exactly marker-wide slot
\newcommand{\annA}{{\ttfamily\color{mkA}[A]}}
\newcommand{\annC}{{\ttfamily\color{mkC}[C]}}
\newcommand{\annP}{{\ttfamily\color{mkP}[P]}}
\newcommand{\annU}{{\ttfamily\color{mkU}[.]}}
\newcommand{\annB}{\hphantom{\ttfamily[A]}}

\lstdefinestyle{annotatedsrc}{
  language=C,
  morekeywords={bool,TRUE,FALSE},
  basicstyle=\fontsize{7}{8}\selectfont\ttfamily,
  keywordstyle=\color{srcKw}\bfseries,
  commentstyle=\color{srcCm}\itshape,
  stringstyle=\color{srcSt},
  columns=fullflexible,
  keepspaces=true,
  showstringspaces=false,
  frame=single,
  rulecolor=\color{black!20},
  framesep=5pt,
  xleftmargin=6pt,
  aboveskip=2pt,
  belowskip=2pt,
  escapeinside={@}{@},
}

\begin{figure}[htbp]
  \centering
  \begin{lstlisting}[style=annotatedsrc]
@\annB\annB@ static void strcpy_url(char *output, const char *url, bool relative)
@\annA\annC@  {
@\annB\annC@    /* we must add this with whitespace-replacing */
@\annB\annC@    bool left = TRUE;
@\annB\annC@    const unsigned char *iptr;
@\annB\annC@    char *optr = output;
@\annB\annC@    const unsigned char *host_sep = (const unsigned char *) url;

@\annB\annC@    if(!relative)
@\annA\annC@      host_sep = (const unsigned char *) find_host_sep(url);

@\annA\annC@    for(iptr = (unsigned char *)url;    /* read from here */
@\annA\annC@        *iptr;         /* until zero byte */
@\annA\annC@        iptr++) {

@\annA\annC@      if(iptr < host_sep) {
@\annA\annC@        *optr++ = *iptr;
@\annB\annC@        continue;
@\annB\annC@      }

@\annA\annC@      switch(*iptr) {
@\annB\annC@      case '?':
@\annA\annC@        left = FALSE;
@\annB\annC@        /* fall through */
@\annB\annC@      default:
@\annA\annP@        if(!ISPRINT(*iptr)) {
@\annA\annU@          snprintf(optr, 4, "%%%02x", *iptr);
@\annB\annU@          optr += 3;
@\annB\annU@        }
@\annB\annB@        else
@\annA\annC@          *optr++=*iptr;
@\annA\annC@        break;
@\annB\annU@      case ' ':
@\annA\annU@        if(left) {
@\annA\annU@          *optr++='%'; /* add a '%' */
@\annB\annU@          *optr++='2'; /* add a '2' */
@\annB\annU@          *optr++='0'; /* add a '0' */
@\annB\annU@        }
@\annB\annB@        else
@\annA\annU@          *optr++='+'; /* add a '+' here */
@\annA\annC@        break;
@\annB\annC@      }
@\annA\annC@    }
@\annA\annC@    *optr = 0; /* zero terminate output buffer */

@\annB\annC@  }
  \end{lstlisting}
  \caption[Coverage-annotated source in the agent workspace]{%
    Excerpt of the coverage-annotated source tree as the agent sees it, here
    \texttt{strcpy\_url} from cURL. The gutter carries an affected marker and
    a coverage marker per line: \annA{} marks membership in the commit-affected
    slice, and \annC, \annP{} and \annU{} denote covered, partially covered, and
    uncovered lines under the current corpus. The agent targets \annA\annU{}
    lines first and \annA\annP{} lines as a fallback. Here the non-printable
    escaping branch and the entire space-encoding branch are unreached, while the
    surrounding loop is covered, which localizes the missing input property to
    URLs containing spaces or non-printable characters.}
  \label{fig:annotated-source}
\end{figure}